%% chapters/behavior_modeling/ch2_rut.tex
%% 第2章　ランダム効用最大化理論の基礎

\section{ランダム効用最大化理論}\label{sec:abms}

本節では行動モデルの理解にあたり，ランダム効用最大化理論（Random Utility Maximization Theory, RUT）を導入する．
RUT はミクロ経済学の消費者行動理論にその源流を持っており，個人の選択行動を効用最大化の枠組みで記述する理論である．
行動モデルの一種であり，
\sref{sec:ch1-network}で整備した記法と，\sref{sec:ch2-uncongested}で扱ったフロー非依存型配分を踏まえ，
本節では\textbf{利用者均衡配分}（User Equilibrium Assignment, UE）を体系的に扱う．
UE は，リンク旅行時間がリンク交通量に依存する混雑ネットワークにおいて，
利用者が自己の旅行時間を最小化しようとする合理的行動の均衡状態を記述するモデルであり，
交通量配分理論の根幹をなす\cite{Wardrop1952,Beckmann1956}．

本節の構成は次のとおりである．
\ssref{subsec:wardrop} で均衡条件を与える Wardrop の第1原則を導入し，
\ssref{subsec:beckmann} でこの条件を等価な凸最適化問題（Beckmann 変換）に変換する．
\ssref{subsec:frank-wolfe} では，この最適化問題の標準的な解法である
\textbf{Frank-Wolfe 法}を他手法との比較も交えて説明する．
\ssref{subsec:braess} では，UE の本質的な性質を浮き彫りにする事例として
\textbf{Braess のパ